% Generated draft: requires mathematical and source review.
\documentclass[11pt]{article}
\usepackage[margin=25mm]{geometry}
\usepackage{fontspec}
\setmainfont{Latin Modern Roman}
\usepackage{amsmath,amssymb,mathtools,unicode-math}
\setmathfont{Latin Modern Math}
\usepackage{enumitem,needspace,xurl,hyperref}
\setlength{\parindent}{0pt}
\setlength{\parskip}{5pt}
\setlength{\emergencystretch}{3em}
\newcommand{\N}{\mathbb{N}}
\newcommand{\Z}{\mathbb{Z}}
\newcommand{\Q}{\mathbb{Q}}
\newcommand{\R}{\mathbb{R}}
\newcommand{\C}{\mathbb{C}}
\newcommand{\F}{\mathbb{F}}
\newcommand{\A}{\mathbb{A}}
\newcommand{\PP}{\mathbb{P}}
\newcommand{\E}{\mathbb{E}}
\newcommand{\eps}{\varepsilon}
\newcommand{\dd}{\,\mathrm{d}}
\begin{document}
\section*{UnsolvedMath ID: 20000728}
\textbf{Source problem code:} catalogue entry
\subsection{Definitions and mathematical statement}
Given a $q$-series $f(q)=\sum_{n\ge0}a_nq^n$ and its weight, the original problem asks for fast heuristic methods to recognize modularity without knowing the level or group. Its suggested directions are Zagier's method using asymptotics at a point, $p$-adic tests for more prime divisibility of coefficients than chance predicts, and comparison of the distribution of $a_p$ with Sato--Tate. These suggestions are not a specified decision algorithm or a uniform probabilistic guarantee.

In the scalar holomorphic congruence setting, the attached research aid asserts a finite-data obstruction. For every even integer $k\ge2$, every $M\ge1$ and every complex prefix $a_0,\ldots,a_M$, there is a holomorphic modular form $F$ of weight $k$ at level
\[
N_k\operatorname{lcm}(1,\ldots,M),\qquad N_2=11,\quad N_4=5,\quad N_k=11\ (k\ge6),
\]
whose expansion agrees through $q^M$. It describes an explicit triangular oldform construction, but does not supply its formula in this summary. For every integer $R>M$, it further asserts that $F(q)+q^R$ is a nonmodular continuation of the same prefix in the same fixed-weight holomorphic congruence interpretation. Thus a prescribed weight and finitely many unrestricted-level coefficients cannot decide modularity; this does not preclude useful heuristics with additional structure.

For normalized Hecke-eigenform candidates the aid separately reports level constraints from local Hecke defects. For example, for a proposed integral weight $k$ and nebentypus $\varepsilon$, the unramified relation is
\[
a_{p^2}=a_p^2-\varepsilon(p)p^{k-1}\qquad(p\nmid N).
\]
A nonzero defect at $p$, under that proposed eigenform interpretation, forces $p\mid N$; distinct such primes therefore give a divisor and lower bound for any candidate level. These tests require the eigenform and character assumptions and are not necessary identities for arbitrary linear combinations of modular forms.

The aid also gives an exact calibration for multiplicativity tests on disjoint pairs of distinct primes. Under a null model in which all involved residues $a_p,a_q,a_{pq}$ are independent and uniform in $\F_\ell$, each test $a_{pq}=a_pa_q$ passes with probability $1/\ell$. If $m$ tests use disjoint coefficient data, their number of passes has distribution $\operatorname{Binomial}(m,1/\ell)$. This is a false-positive calculation for the stated residue model, not a probability that an arbitrary input series is modular. Appropriate eigenform, normalization and distributional hypotheses would also be needed to turn the original Sato--Tate suggestion into a test.
\subsection{Short English statement}
Find practical weight-only modularity heuristics. The supplied aid shows that finite coefficients cannot decide modularity when the level is unrestricted, but retains conditional eigenform tests that can constrain a candidate level and quantify random residue coincidences.
\subsection{Sources}
\begin{itemize}
\item Local UnsolvedMath registry, permanent ID $20000728$, source code catalogue entry; canonical \texttt{aim-analytic-number-theory-notes.json}, \texttt{notes[91]}.
\item AIM workshop ``Mock modular forms in combinatorics and arithmetic geometry,'' March 8--12, 2010, problem 1.1 in the two-page list \emph{Mock Modular Forms}, edited by Sharon Anne Garthwaite: \url{https://aimath.org/WWN/mockmodular/mockmodular.pdf}. The supplied verification field reports direct PDF inspection; this draft relies on that supplied record. The original remarks name Zagier and suggest $p$-adic divisibility and Sato--Tate comparisons; no further bibliography is supplied.
\item Attached research attempt 1: partial result, recorded run status open, rigor proved, novelty confidence low; registry status partially solved, classification PARTIAL-PROGRESS. The machine-generated aid is not peer reviewed and requires independent verification before citation.
\end{itemize}
\end{document}
